\documentclass{article}
\usepackage[table]{xcolor}  %\cellcolor[gray]{0.9}  \cellcolor{red}
\usepackage{multirow}
%\usepackage{siunitx}
 
 \begin{document}
\title{Folding a pattern}
\author{Alejandro Rivero}\maketitle
\begin{abstract}
 We propose a reorganisation of the standard model and their mesons in order
 to build supersymmetric multiplets. The presentation is open to improvements
 to choose the adequate candidates in each recombination.
\end{abstract}



%\centering


%the folded pattern here



I have told elsewhere how a chain of four Koide equations does a good work to predict SM mass.
Gray cells are inputs
All triplets in sequence meet Koide equation. 
For $(bcs)$, sign of $\sqrt {m_s}$ is minus.

\begin{center}
\begin{tabular}{r|r||c|c||}
         &  {pdg 2012}       & {exact} & {rotated}  \\
\hline
t & $173.5 \pm 1.0 $ & \cellcolor[gray]{0.9} 174.10 & 173.26     \\
b & $4.18 \pm 0.03$ & 3.64 &4.197    \\
c & $1.275 \pm 0.025$  & 1.698  &1.359   \\
$\tau$ & $1.77682 (16)$ & 1.698 &1.776968  \\
s & $95 \pm 5$ & 121.95 & 92.275           \\
$\mu$ &105.65837  & 121.95 & \cellcolor[gray]{0.9}105.6584       \\
d &$\sim 4.8$ & 8.75 &  5.32    \\
u &$\sim 2.3$ & \cellcolor[gray]{0.9} 0 & .03564          \\
e &0.5109989 & 0 & \cellcolor[gray]{0.9} .5109989       \\
\hline
\end{tabular}
\end{center}

Lets exploit this identification of some levels to try to fold them into a susy-like configuration.

First step, the SM as it is. For each QCD string, consider its fundamental state.

$$
\begin{array}{lllllll} 
&\nu_1, t_{rgb}&  &   & & \\ 
&\nu_2, b_{rgb}& 
      B^+,B_c^+ & 
      {bu}, {bc}
      & {bb}, {bs}, {bd} & 
      \eta_b, B^0, B^0_s,  \bar B^0, \bar B^0_s \\ 
\stackrel{\bar c\bar c}{cc},\stackrel{\bar c\bar u}{cu}&
\tau, c_{rgb} &
       D^+, D_s^+&
       sc,dc  &
       & 
       \eta_c, D^0, \bar {D^0}\\ 
\stackrel{\bar u\bar u}{uu}&
\mu, s_{rgb} &
         \pi^+, K^+& 
         su, du&
         ss, sd, dd &
         \eta_8, \pi^0, K^0, \bar K^0 \\ 
&\nu_3,  d_{rgb} \\  
&e, u_{rgb}\end{array}
$$

For charged particles, the antiparticle is in the same level. I ommit the ``+ antiparticle'' remark

Now, second step, we fold t and u

$$
\begin{array}{lllllll} 
&\nu_2, b_{rgb}, e, u_{rgb}& 
      B^+,B_c^+ & 
      {bu}, {bc}
      & {bb}, {bs}, {bd} & 
      \eta_b, B^0, B^0_s,  \bar B^0, \bar B^0_s \\ 
\stackrel{\bar c\bar c}{cc},\stackrel{\bar c\bar u}{cu}&
\tau, c_{rgb} &
       D^+, D_s^+&
       sc,dc  &
       & 
       \eta_c, D^0, \bar {D^0}\\ 
\stackrel{\bar u\bar u}{uu}&
\mu, s_{rgb} , \nu_1, t_{rgb}&
         \pi^+, K^+& 
         su, du&
         ss, sd, dd &
         \eta_8, \pi^0, K^0, \bar K^0 \\ 
&\nu_3,  d_{rgb} 
\end{array}
$$

Next steps are arbitrary, we could fold any combination of up-like quarks

$$
\begin{array}{lllllll} 
&\nu_2, b_{rgb}, e, u_{rgb}& 
      B^+,B_c^+ & 
      {bu}, {bc}
      &  {bs}, {bd} & 
      \eta_b, B^0, B^0_s,  \bar B^0, \bar B^0_s \\ 
\stackrel{\bar c\bar c}{cc},\stackrel{\bar c\bar u}{cu}&
\tau, c_{rgb} , \nu_3,  d_{rgb}&
       D^+, D_s^+&
       sc,dc  & 
       bb,dd &
       \eta_c, D^0, \bar {D^0}\\ 
\stackrel{\bar u\bar u}{uu}&
\mu, s_{rgb} , \nu_1, t_{rgb}&
         \pi^+, K^+& 
         su, du&
         ss, sd&
         \eta_8, \pi^0, K^0, \bar K^0 
\end{array}
$$

And last, and even more arbitrary, is to decide which neutral meson
to fold into the intermediate level. Just choose one

$$
\begin{array}{lllllll} 
&\nu_2, b_{rgb}, e, u_{rgb}& 
      B^+,B_c^+ & 
      {bu}, {bc}
      &  {bs}, {bd} & 
       B^0, B^0_s,  \bar B^0, \bar B^0_s \\ 
\stackrel{\bar c\bar c}{cc},\stackrel{\bar c\bar u}{cu}&
\tau, c_{rgb} , \nu_3,  d_{rgb}&
       D^+, D_s^+&
       sc,dc  & 
       bb,dd &
       \eta_b, \eta_c, D^0, \bar {D^0}\\ 
\stackrel{\bar u\bar u}{uu}&
\mu, s_{rgb} , \nu_1, t_{rgb}&
         \pi^+, K^+& 
         su, du&
         ss, sd&
         \eta_8, \pi^0, K^0, \bar K^0 
\end{array}
$$

The $\pm 4/3$ diquarks, at last, I guess they are related to 
electroweak symmetry breaking via some condensation. They could stay
as there are, or perhaps the charm should be at an upper level

$$
\begin{array}{lllllll} 
\stackrel{\bar c\bar c}{cc}&
\nu_2, b_{rgb}, e, u_{rgb}& 
      B^+,B_c^+ & 
      {bu}, {bc}
      &  {bs}, {bd} & 
       B^0, B^0_s,  \bar B^0, \bar B^0_s \\ 
\stackrel{\bar c\bar u}{cu}&
\tau, c_{rgb} , \nu_3,  d_{rgb}&
       D^+, D_s^+&
       sc,dc  & 
       bb,dd &
       \eta_b, \eta_c, D^0, \bar {D^0}\\ 
\stackrel{\bar u\bar u}{uu}&
\mu, s_{rgb} , \nu_1, t_{rgb}&
         \pi^+, K^+& 
         su, du&
         ss, sd&
         \eta_8, \pi^0, K^0, \bar K^0 
\end{array}
$$

It could be. Probably all the ambiguities in the above folding,
including this last one, are related.

Final, let me prettify it by adding the antiparticles. Note that we have the 
required quantity of bosons for three generations of the SM, but that
the not-SM charged bosons (the uu etc) can not be arranged in a
pattern of three generations of Dirac particles; it is because of it
that we can guess they have a different role that being
a plain sfermion.
 
$$
\begin{array}{||l|l|llll||} 
\hline
\stackrel{\bar c\bar c}{cc}&
\nu_2, b_{rgb}, e, u_{rgb}& 
      B^\pm,B_c^\pm & 
      \stackrel{\bar b\bar u}{bu}, \stackrel{\bar b\bar c}{bc}
      &  \stackrel{\bar b \bar s}{bs}, \stackrel{\bar b\bar s}{bd} & 
       B^0, B^0_s,  \bar B^0, \bar B^0_s \\ 
\stackrel{\bar c\bar u}{cu}&
\tau, c_{rgb} , \nu_3,  d_{rgb}&
       D^\pm, D_s^\pm&
       \stackrel{\bar s\bar c}{sc},\stackrel{\bar d\bar c}{dc}  & 
       \stackrel{\bar b\bar b}{bb},\stackrel{\bar d\bar d}{dd} &
       \eta_b, \eta_c, D^0, \bar {D^0}\\ 
\stackrel{\bar u\bar u}{uu}&
\mu, s_{rgb} , \nu_1, t_{rgb}&
         \pi^\pm, K^\pm& 
         \stackrel{\bar s\bar u}{su}, \stackrel{\bar d\bar u}{du}&
         \stackrel{\bar s\bar s}{ss}, \stackrel{\bar s\bar d}{sd}&
         \eta_8, \pi^0, K^0, \bar K^0 \\
\hline
\end{array}
$$

Remember that the point of the grouping in fermion side is that we were
looking at the koide triplets of the waterfall; any of the triplets
is build by taking a fermion for each of the three lines.

A lateral question is, can be build mass values so that any of the 
eight possible triplets is solution?

Let me consider a simpler case, that of, say, $m_u = 0$. Call
$$k=k^+= 2 + {\sqrt 3}$$
$$k^-= 2 - {\sqrt 3}$$
$$k^+k^-=1$$
and check that $(k, 8, k^-)$ and $(k, 0, k^-)$ are Koide
triplets. We can use them to build ``mass spectra'' that saturate
all the equations.

$$ \begin{array}{|c|c|c|}
\begin{array}{ll}
 0     & \\
 1     & \\
 k^2  &k^2 
\end{array} &

\begin{array}{ll}
 0      &           \\
 1      &            \\
 k^{-2} & k^{-2}
\end{array} &


\begin{array}{ll}
 0   &        \\
 1   &         \\
 k^2 & k^{-2}
\end{array} \\


\begin{array}{|ll|}
 0     & \\
 1     & 1\\
 k^2  &k^2 
\end{array} 


\begin{array}{|ll|}
 0     & \\
 1     & k^4\\
 k^2  &k^2 
\end{array} 


&

\begin{array}{|ll|}
 0      &           \\
 1      & 1          \\
 k^{-2} & k^{-2}
\end{array}

\begin{array}{|ll|}
 0      &           \\
 1      & k^{-4}          \\
 k^{-2} & k^{-2}
\end{array} &


\begin{array}{|ll|}
 0   &        \\
 1   &   1      \\
 k^2 & k^{-2}
\end{array} \\



%%%%%%%%%%%%%%%%%%%%%%%%%%%%%


\begin{array}{|ll|}
 0     & 0\\
 1     & 1\\
 k^2  &k^2 
\end{array} 

\begin{array}{|ll|}
 0     & 8k\\
 1     & 1\\
 k^2  &k^2 
\end{array} 


\begin{array}{|ll|}
 0     & 0\\
 1     & k^4\\
 k^2  &k^2 
\end{array} 


&

\begin{array}{|ll|}
 0      & 0        \\
 1      & 1          \\
 k^{-2} & k^{-2}
\end{array}

\begin{array}{|ll|}
 0      & 8k^-        \\
 1      & 1          \\
 k^{-2} & k^{-2}
\end{array}

\begin{array}{|ll|}
 0      &   0        \\
 1      & k^{-4}          \\
 k^{-2} & k^{-2}
\end{array} &


\begin{array}{|ll|}
 0   &    0    \\
 1   &   1      \\
 k^2 & k^{-2}
\end{array} \\

%%%%%%%%%%%%%%%%%%%%%%%%%%%%%
\hline
\hline

\begin{array}{|ll|}
 0   & 0\\
 k^- & k^-\\
 k   &k 
\end{array} 

\begin{array}{|ll|}
 0     & 8\\
 k^-     & k^-\\
 k  &k 
\end{array} 


\begin{array}{|ll|}
 0     & 0\\
 k^{-2}     & k^2\\
 1  &1 
\end{array} 


&

\begin{array}{|ll|}
 0      & 0        \\
 k      & k          \\
 k^{-} & k^{-}
\end{array}

\begin{array}{|ll|}
 0      & 8        \\
 k      & k          \\
 k^{-} & k^{-}
\end{array}

\begin{array}{|ll|}
 0      &   0        \\
 k^2      & k^{-2}          \\
 1 & 1
\end{array} &


\begin{array}{|ll|}
 0   &    0    \\
 1   &   1      \\
 k^2 & k^{-2}
\end{array} \\
\hline
\hline
\end{array}
$$

In this case, we see that there is really only three possibilites: we can 
keep full degenerated with the trivial triplet $(0, 2 - \sqrt 3 , 2+\sqrt 3 )$ or
we can do one of two possible mutations: either change one zero to
a nonzero from the complementary triplets (8 or 96, thus), or to multiply 
one of the $2 - \sqrt 3$ masses with $k^4$, scaling it to $k^3=26+ 15 \sqrt 3$. 

Phenom-wise, the two changes are interesting. In the former, we could think
that the broken double is the $u,b$ pair, and then the $b$ quark gains -or keeps-
a mass above the other two doubles. In the later, it should be the
$t,s$ pair, separating the top quark still not far enough, but the most
allowed by the set of eight simultaneus Koide equations (at least,
the most with a massless quark still there; we should work now the
general case).

$$
\begin{array}{c}
\hline
\begin{array}{||ll||}
 0     & 96\\
 k^-     & k^-\\
 k  &k 
\end{array} 
\begin{array}{ll}
 0     & 8\\
 k^-     & k^-\\
 k  &k 
\end{array} 
\begin{array}{||ll}
 0      & 0        \\
 k^{-} & k^{-} \\
 k      & k          
\end{array}
\begin{array}{|||ll||}
 0   &    0    \\
 k^3 & k^{-} \\ 
 k   &   k     
\end{array} \\
\hline
\end{array}
\dots
\begin{array}{|ll|}
 96     & 8\\
 k^-     & k^-\\
 k  &k 
\end{array} 
\dots
$$



If we put some mass scale in GeV, the likeliness with the SM is apparent

$$
\begin{array}{c}
\hline
\begin{array}{||ll||}
 0  & 43.6918 \\ 
 0.12195 & 0.12195 \\
1.69854 & 1.69854
\end{array} 
\begin{array}{ll}
 0     & 3.64098 \\
 0.12195 & 0.12195 \\
 1.69854 & 1.69854
\end{array} 
\begin{array}{||ll}
 0      & 0        \\
 0.12195 & 0.12195 \\
 1.69854 & 1.69854 \\
\end{array}
\begin{array}{|||ll||}
 0   &    0    \\
 1.69854  & 0.008755 \\
 0.12195 & 0.12195
\end{array} \\
\hline
\end{array}
$$



Ok; if we are going to use these sextets as a startpoint to look
for continous families of non-zero solutions, we can also change 
both zeros to a still degenerated pair (8,8) or
(96,96) or to the non-degenerated (8,96). But here we see that
the general case of two degenerated pairs is uninteresting: 
there is always the possibility of choosing from up to four 
different solutions for a pair of masses. In any case, we could
look for accidental degeneracies.

So a more interesting question 
is if there are solutions with zero or one degenerated pair at most.
A numerical exploration, e.g with {\tt octave} finds at least four candidates,
they are uninteresting, if they are expected to have some hint about
the top quark:
                      
%[1.972064777327935,12.1354051054384]
%[.2471462221416923,1.972064777327935]
%[.2471462221416923,4.974009900990099]
%[4.974009900990099,12.1354051054384]           3?           
%                      
                      
%[.2626126990216766,5.046782114241406]          
%[3.752061572292468,5.046782114241406]          2?                    
%[.003668609282629659,.2626126990216766]
%[.003668609282629659,3.752061572292468]    
                        
%[7.336050724637682,21.41397288842544]          4?
%[4.38487023743788,21.41397288842544]
%[3.930230298318919,7.336050724637682]
%[3.930230298318919,4.38487023743788]

%[.002862920633477994,.2721622021607222]
%[.002862920633477994,3.747673216132368]
%[3.747673216132368,5.091511936339523]          1?
%[.2721622021607222,5.091511936339523]
\begin{center}
\begin{tabular}{|ll|ll|ll|ll|}
0.000055046 &  174.1  & 0.000091997  & 174.1  & 0.0722 & 174.1& 5.864 & 174.1\\
 0.49746& 94.32  & 0.4714 & 96.229  & 4.597 &  29.24 & 7.299 & 20.432\\
6.7159 & 6.7159 & 6.8354 & 6.8354 & 1.1821 & 1.1821 & 0.3796& 0.3796 \\
\hline
\end{tabular}
\end{center}

 
 They are not directly connected to the zeroed solutions.
 The next approach is to use Mathematica or Maxima to find
 solutions perhaps missed in the numerical exploration
  This is explained in
 
 {\tt http://www.physicsforums.com/showthread.php?t=551549\&page=7}
 
 {\tt http://www.physicsforums.com/showpost.php?p=4270855\&postcount=100}
 
 and what happens is very peculiar; the polynomials in the resolvent
 produce the above four solutions but each of them
 have extra real triplets that can be used to build zero-less but degenerated
 solutions; these solutions are really inside the continous spectra
 of answers, but they are specially signaled in the resolvent and they
 are also connected in some way to the above zeroed solutions. In
 particlar, the strange-charm-bottom triplet
 $$
[2-\sqrt 3,1,2\sqrt 3 -2]   
$$
and the charm-bottom-top
$$
[1, 2\sqrt 3-2,7\sqrt 3-2] 
$$
solutions appear explicitly. So we have on one side a zeroed solution, and
on other side, from the resolvent, a particular zeroless solution, that happens
to be related to our folding:

$$
\begin{array}{|ll|}
\hline
 3.64098 & 0 \\
 1.69854 & 1.69854 \\
 0.12195 & 0.12195 \\
\hline
\end{array}
\dots
\begin{array}{|ll|}
\hline
 b  & u \\
 d  & c \\
 s  & t \\
\hline
\end{array} 
\dots
\begin{array}{|ll|}
\hline
 3.640 & 3.640 \\
 1.698 & 1.698 \\
 0.1219 & 174.1 \\
\hline
\end{array}
$$
 
Note that some of the zeroed solutions were also able to produce good values for the down
quark. Perhaps the right symmetry group is beyond $S_4$

\end{document}
